% Source: https://artofproblemsolving.com/community/c3h2171675_confused_triangle
% Stable identity: aops-topic-2171675-q01
Let $K=[ABC]$.

\begin{asy}
usepackage("amsmath");
size(7cm);
defaultpen(fontsize(10pt));
pair A=(1.5,7),B=(0,0),C=(10,0);
real r1=2/12, r2=3/12, r3=7/12;
pair G=B+r1*(C-B), F=B+(r1+r3)*(C-B);
pair O=extension(G,G+(A-B),F,F+(A-C));
pair M=extension(O,O+(C-B),A,B), Np=extension(O,O+(C-B),A,C);
pair Ep=extension(O,O+(C-A),A,B), H=extension(O,O+(B-A),A,C);
fill(M--O--Ep--cycle,rgb(1,0.85,0.85)); fill(Np--O--H--cycle,rgb(0.85,1,0.85)); fill(G--O--F--cycle,rgb(0.85,0.9,1));
draw(A--B--C--cycle,linewidth(0.9));
draw(M--Np,blue); draw(Ep--F,red); draw(G--H,deepgreen);
dot(A);dot(B);dot(C);dot(O);dot(M);dot(Np);dot(Ep);dot(H);dot(G);dot(F);
label("$A$",A,N);label("$B$",B,SW);label("$C$",C,SE);label("$O$",O,1.3*dir(109),fontsize(9pt));
label("$M$",M,W);label("$E$",Ep,W);label("$N$",Np,NE);label("$H$",H,NE);label("$G$",G,S);label("$F$",F,S);
label("$4$",(M+O+Ep)/3,fontsize(8pt));label("$9$",(Np+O+H)/3,fontsize(8pt));label("$49$",(G+O+F)/3,fontsize(9pt));
label("\mbox{}",(min(currentpicture,user=true)+max(currentpicture,user=true))/2,black); // keeps the text after the figure black when LaTeX typesets the labels inline
\end{asy}

\textbf{Three similar triangles.} Each small triangle has its sides parallel to the sides of $ABC$: for example, triangle $MOE$ has $ME$ on $AB$, $MO\parallel BC$ and $OE\parallel AC$. So each is similar to $ABC$, and since areas scale with the square of the ratio, the ratios are
\[r_1=\sqrt{\dfrac{4}{K}}\ (\triangle MOE),\qquad r_2=\sqrt{\dfrac{9}{K}}\ (\triangle NOH),\qquad r_3=\sqrt{\dfrac{49}{K}}\ (\triangle GOF).\]
In particular $MO=r_1\cdot BC$, $ON=r_2\cdot BC$ and $GF=r_3\cdot BC$ (these are the sides parallel to, or on, $BC$).

\textbf{Parallelograms.} $BMOG$ is a parallelogram ($MO\parallel BG$ and $OG\parallel MB$), so $BG=MO$. Likewise $OFCN$ is a parallelogram ($ON\parallel FC$ and $OF\parallel NC$), so $FC=ON$. Since $BG+FC=MO+ON=MN<BC$, the points lie in the order $B,G,F,C$ on side $BC$, and
\[BC=BG+GF+FC=\left(r_1+r_3+r_2\right)BC .\]
Hence $r_1+r_2+r_3=1$, that is,
\[\dfrac{2+3+7}{\sqrt K}=1,\qquad\sqrt K=12,\qquad K=\boxed{144}.\]
